\documentclass[11pt]{article}
\usepackage[a4paper,margin=1in]{geometry}
\usepackage{fontspec}
\usepackage[boldfont=false]{xeCJK}
\setCJKmainfont{FandolSong-Regular.otf}
\usepackage{amsmath}
\usepackage{amssymb}
\usepackage{array}
\usepackage{booktabs}
\usepackage{graphicx}
\usepackage{adjustbox}
\usepackage{enumitem}
\setlist[itemize]{topsep=2pt,partopsep=0pt,parsep=0pt,itemsep=2pt,leftmargin=1.4em}
\setlist[enumerate]{topsep=2pt,partopsep=0pt,parsep=0pt,itemsep=2pt,leftmargin=1.6em}
\usepackage{parskip}
\usepackage{fvextra}
\fvset{breaklines=true,breakanywhere=true,fontsize=\small}
\setlength{\parindent}{0pt}

\begin{document}

\begin{center}
  {\LARGE\bfseries Chemical reactions}\\[0.35em]
  {\large Igcse Chemistry}
\end{center}
\vspace{0.6em}

\section*{Physical and chemical changes}

A \textbf{physical change} 物理变化 does not make a new substance. The substance only changes its state or shape, and the change can usually be reversed. Melting ice and dissolving sugar are physical changes.

A \textbf{chemical change} 化学变化 (a chemical reaction) makes one or more new substances and is usually hard to reverse. Signs of a chemical change include a colour change, a gas being given off, an energy change, or a \textbf{precipitate} 沉淀 (a solid) forming.

\begin{center}
\adjustbox{max width=\linewidth}{%
\begin{tabular}{ll}
\toprule
\textbf{Physical change} & \textbf{Chemical change} \\
\midrule
no new substance made & new substance(s) made \\
easy to reverse & usually hard to reverse \\
e.g. melting, boiling, dissolving & e.g. burning, rusting \\
\bottomrule
\end{tabular}}
\end{center}

\section*{Rate of reaction}

The \textbf{rate of reaction} 反应速率 tells you how fast the \textbf{reactants} 反应物 change into \textbf{products} 生成物.

\subsection*{Collision theory}

\textbf{Collision theory} 碰撞理论 explains what controls the rate. For a reaction to happen, the \textbf{particles} 粒子 must \textbf{collide} 碰撞, and they must collide with enough energy. The least energy they need is the \textbf{activation energy} 活化能 ($E_a$). A faster rate happens when the particles have successful collisions more often.

\subsection*{What changes the rate}

\begin{center}
\adjustbox{max width=\linewidth}{%
\begin{tabular}{lll}
\toprule
\textbf{Change} & \textbf{Effect} & \textbf{Reason (collision theory)} \\
\midrule
increase \textbf{concentration} 浓度 of a solution & faster & more particles in the same volume, so collisions happen more often \\
increase \textbf{pressure} 压强 of gases & faster & particles are pushed closer, so collisions happen more often \\
increase \textbf{surface area} 表面积 of a solid & faster & more particles are exposed, so collisions happen more often \\
increase temperature & faster & particles gain \textbf{kinetic energy} 动能 and move faster, so collisions are more frequent and more of them have enough energy \\
add a \textbf{catalyst} 催化剂 & faster & the catalyst lowers the activation energy, so more collisions are successful \\
\bottomrule
\end{tabular}}
\end{center}

\subsection*{Catalysts}

A catalyst speeds up a reaction but is not used up — it is unchanged at the end. It works by lowering the activation energy. \textbf{Enzymes} 酶 are biological catalysts (catalysts that work in living things).

\subsection*{Measuring the rate}

You can follow a reaction over time in these ways:

\begin{itemize}
\item \textbf{Change in mass}: stand the flask on a balance. If a gas escapes, the mass falls. Record the mass at regular times.

\item \textbf{Volume of gas}: collect the gas in a gas \textbf{syringe} 注射器 and read its volume at regular times.

\item \textbf{Formation of a precipitate}: in a reaction that turns cloudy, time how long it takes for a mark under the flask to disappear.

\end{itemize}

On a graph of product against time, the line is \textbf{steepest at the start} (fastest rate), becomes less steep as reactants are used up, and goes \textbf{flat} when the reaction has finished.

\section*{Reversible reactions and equilibrium}

Some reactions are \textbf{reversible} 可逆反应: the products can react to form the reactants again. We show this with the symbol $\rightleftharpoons$.

\subsection*{Changing the direction}

A clear example uses \textbf{hydrated} 水合 and \textbf{anhydrous} 无水 compounds:

\begin{itemize}
\item Blue hydrated copper(II) sulfate, when heated, loses its water to become white anhydrous copper(II) sulfate. Adding water turns it blue again.

\item Pink hydrated cobalt(II) chloride loses water when heated to become blue anhydrous cobalt(II) chloride. Adding water turns it pink again.

\end{itemize}

Adding water to the anhydrous solid (and seeing the colour return) is used as a test for water.

\subsection*{Equilibrium}

In a \textbf{closed system} 密闭系统 (where nothing enters or leaves), a reversible reaction reaches \textbf{equilibrium} 平衡 when:

\begin{itemize}
\item the rate of the \textbf{forward reaction} 正反应 equals the rate of the \textbf{reverse reaction} 逆反应, and

\item the concentrations of the reactants and products are no longer changing.

\end{itemize}

\subsection*{Changing the position of equilibrium}

The \textbf{position of equilibrium} 平衡位置 tells you whether there are more reactants or more products. You can move it:

\begin{itemize}
\item \textbf{Temperature}: heating moves the equilibrium in the direction that takes in heat (the \textbf{endothermic} 吸热反应 direction); cooling moves it in the direction that gives out heat (the \textbf{exothermic} 放热反应 direction).

\item \textbf{Pressure} (gases): more pressure moves the equilibrium to the side with fewer gas molecules.

\item \textbf{Concentration}: adding more of a substance moves the equilibrium to the other side, to use it up.

\item A \textbf{catalyst} does \textbf{not} move the position of equilibrium. It only helps the reaction reach equilibrium faster.

\end{itemize}

\subsection*{The Haber process}

The \textbf{Haber process} 哈伯法 makes ammonia:

\[
\text{N}_2(g) + 3\text{H}_2(g) \rightleftharpoons 2\text{NH}_3(g)
\]

\begin{itemize}
\item The \textbf{hydrogen} 氢气 comes from \textbf{methane} 甲烷 (natural gas); the \textbf{nitrogen} 氮气 comes from the air.

\item Typical conditions: a temperature of $450\,^{\circ}\text{C}$, a pressure of about $200$ atm ($20\,000$ kPa), and an \textbf{iron} 铁 catalyst.

\end{itemize}

\subsection*{The Contact process}

The \textbf{Contact process} 接触法 turns sulfur dioxide into sulfur trioxide:

\[
2\text{SO}_2(g) + \text{O}_2(g) \rightleftharpoons 2\text{SO}_3(g)
\]

\begin{itemize}
\item The sulfur dioxide comes from burning \textbf{sulfur} 硫 (or roasting sulfide ores); the \textbf{oxygen} 氧气 comes from the air.

\item Typical conditions: a temperature of $450\,^{\circ}\text{C}$, a pressure of about $2$ atm ($200$ kPa), and a \textbf{vanadium(V) oxide} 五氧化二钒 catalyst.

\end{itemize}

\subsection*{Why these conditions are chosen}

The conditions are a \textbf{compromise} 折中:

\begin{itemize}
\item A higher pressure would give more product, but very high pressure is dangerous and expensive, so a medium pressure is used.

\item A lower temperature would give more product (both forward reactions are exothermic), but the reaction would be too slow, so a fairly high temperature is used to keep a good rate.

\item The catalyst speeds up the reaction without changing the position of equilibrium, which lowers cost.

\end{itemize}

\section*{Redox}

\textbf{Oxidation} 氧化 and \textbf{reduction} 还原 always happen at the same time, in what is called a \textbf{redox reaction} 氧化还原反应. ('Redox' is short for reduction–oxidation.)

There are three ways to describe oxidation and reduction:

\begin{itemize}
\item \textbf{Oxygen}: oxidation is the gain of oxygen; reduction is the loss of oxygen.

\item \textbf{Electrons} 电子: oxidation is the loss of electrons; reduction is the gain of electrons.

\item \textbf{Oxidation number} 氧化数: in oxidation the oxidation number goes up; in reduction it goes down.

\end{itemize}

A useful memory aid is \textbf{OIL RIG}: \emph{Oxidation Is Loss, Reduction Is Gain} — of electrons.

\subsection*{Oxidation number rules}

The oxidation number is shown by a Roman numeral, as in iron(II) and iron(III). The rules are:

\begin{itemize}
\item An element that is not combined has an oxidation number of $0$.

\item A single-atom ion has an oxidation number equal to its charge (so $\text{Na}^{+}$ is $+1$).

\item The oxidation numbers in a compound add up to $0$.

\item The oxidation numbers in an ion add up to the charge on the ion.

\end{itemize}

\subsection*{Oxidising and reducing agents}

\begin{itemize}
\item An \textbf{oxidising agent} 氧化剂 oxidises another substance, and is itself reduced.

\item A \textbf{reducing agent} 还原剂 reduces another substance, and is itself oxidised.

\end{itemize}

Some redox reactions show clear colour changes:

\begin{itemize}
\item Acidified \textbf{potassium manganate(VII)} 高锰酸钾 is purple. When it acts as an oxidising agent it is reduced, and the purple colour fades to colourless.

\item \textbf{Potassium iodide} 碘化钾 is colourless. When it is oxidised, red-brown \textbf{iodine} 碘 is formed.

\end{itemize}


\end{document}
